\section{Introduction}
\label{sec:intro}

Large Language Models (LLMs) are increasingly deployed in domains requiring spatial intelligence, from user interface automation to robotic planning. A critical capability for these applications is the ability to learn and reason about novel spatial configurations presented in-context. For instance, a model might be given the coordinate definitions of a new tool or a specific room layout and must subsequently interact with named components of that geometry.

The significance of this research lies in understanding how symbolic representations of geometry, such as \svg coordinate strings, are internalized during in-context learning. Recent studies have demonstrated that LLMs can learn the topology of hidden graphs or 2D lattices through repeated examples~\cite{tang2026concept, xiong2026lattice}. However, a key question remains: does the ability to map a verbal label to a coordinate—which we call \textbf{\addressing}—enable the model to "perceive" the spatial relationships between those coordinates?

Despite the impressive ability of models like \gpt to retrieve exact coordinate strings, there is a clear gap in their ability to perform genuine spatial reasoning. On one hand, the model can act as a perfect lookup table for geometric definitions. On the other hand, it often fails at simple queries like "Is part A above part B?", even when it has just correctly retrieved the coordinates for both parts. This suggests that \addressing and spatial perception may be decoupled in the model's latent space.

In this paper, we examine this decoupling by evaluating \gpt on the \kilogram dataset~\cite{ji2022kilogram}, which contains thousands of unique 2D tangram shapes defined by \svg coordinates. We vary the number of examples ($N$ shots) to identify the "addressability threshold"—the point where symbolic mapping becomes stable—and test whether reaching this threshold triggers a phase transition in \reasoning performance.

Our quantitative evaluation shows that \gpt achieves {\bf 100\% \addressing accuracy} in zero-shot settings, requiring no repetition to map names to complex coordinate strings. However, \reasoning accuracy starts at a mere {\bf 60\%} with one shot and only reaches {\bf 80\%} after 20 shots, showing no sharp phase transition. These findings indicate that while LLMs are exceptional at symbolic retrieval, they lack an intrinsic spatial "eye" for novel geometries learned in-context.

We make the following contributions:
\begin{itemize}[leftmargin=*,itemsep=0pt,topsep=0pt]
    \item We define and distinguish between {\bf \addressing} and {\bf \reasoning} in the context of in-context learning of 2D geometries.
    \item We demonstrate that \gpt exhibits perfect \addressing for complex \svg shapes without requiring repetition, challenging the necessity of the "phase transitions" observed in simpler graph-learning tasks.
    \item We provide empirical evidence for the decoupling of symbolic lookup and spatial perception, showing that \reasoning performance scales much more slowly than \addressing.
    \item We evaluate the impact of \cotask and novel nomenclature, finding that neither significantly bridges the gap between retrieval and reasoning.
\end{itemize}
